\chapter*{Словарь терминов}             % Заголовок
%\addcontentsline{toc}{chapter}{Словарь терминов}  % Добавляем его в оглавление
%
\textbf{TDP} : модели, ценообразование в которых зависит от времени (англ. ``time-dependent models'')\\

\textbf{SDP} : модели, ценообразование в которых зависит от состояния экономики (англ. ``state-dependent models'')\\

\textbf{DSGE} : динамические стохастические модели общего равновесия\\

\textbf{CPI} : индекс потребительских цен, статистический показатель, измеряющий изменения уровня цен на товары и услуги, потребляемые домохозяйствами\\

\textbf{Menu Costs} : издержки, связанные с изменением цен фирмами, включая затраты на перераспределение информации, изменение ценников, уведомление клиентов и т. д.\\

\textbf{Inflation Targeting} : стратегия монетарной политики, направленная на поддержание инфляции на заданном уровне или в пределах определенного диапазона\\

\textbf{Sticky Prices} : жесткость цен, которая возникает из-за различных факторов, препятствующих быстрому реагированию цен на экономические шоки\\

\textbf{Real Rigidities} : реальные жесткости, связанные с ограничениями в экономике, такими как издержки меню и другие трения, которые влияют на гибкость цен\\

\textbf{Nominal Rigidities} : номинальные жесткости, связанные с неготовностью фирм изменять цены из-за факторов, таких как контракты и ожидания\\

\textbf{Phillips Curve} : кривая Филлипса, отражающая обратную зависимость между уровнем безработицы и уровнем инфляции


%\textbf{TeX} : Cистема компьютерной вёрстки, разработанная американским профессором информатики Дональдом Кнутом
%
%\textbf{панграмма} : Короткий текст, использующий все или почти все буквы алфавита
